\chapter{Deconvolution as a damped Weil criterion}\label{chap:dc}



\paragraph{What is not claimed.} The half plane $\re s>\frac78$ has not been refereed. The eigenvalues of the ladder are variational upper bounds, the converged values are extrapolations of such bounds with an error estimate, and none of them is a certificate.

\paragraph{The question.}\label{dc:sec:question} Let $\xi$ have zeros $\rho=\beta+i\gamma$, let $\gamma_\rho=-i(\rho-\frac12)=\gamma-i(\beta-\frac12)$ be the zeros of $\Xi(\theta)=\xi(\frac12+i\theta)$, and put $\sigma_{\max}=\sup_\rho\re\rho\in[\frac12,1]$. We say that a level $\varepsilon>0$ is \emph{to the right of the zeros} if $\beta<\frac12+\varepsilon$ for every zero. Every $\varepsilon\ge\frac12$ is, since $\beta<1$, whether or not $\sigma_{\max}<1$; the latter is the quasi Riemann hypothesis. For real $\alpha$ write $v_\alpha(\theta)=\re\frac{\xi'}{\xi}(\alpha+i\theta)$ and $b_\alpha=v_\alpha(0)=\frac{\xi'}{\xi}(\alpha)$, so that $\rho_\varepsilon=\pi^{-1}v_{1/2+\varepsilon}$ is the density \eqref{eq:rho}: for $\varepsilon>0$ the Poisson smoothing of the zeros (Proposition~\ref{prop:struct}) and, when $\varepsilon$ is to the right of the zeros, the Thorin density of the shifted clock at basepoint $\frac12+\varepsilon$ (Theorem~\ref{thm:poisson}). Write $P_a(x)=a/(\pi(a^2+x^2))$ for the Cauchy density of scale $a>0$. Theorem~\ref{tf:thm:deconv} of Chapter~\ref{chap:tf} says that, for $\varepsilon_0$ to the right of the zeros and $0\le\varepsilon<\varepsilon_0$, there is a positive measure $\nu$ with $\int(1+\theta^2)^{-1}\nu(d\theta)<\infty$ and $\rho_{\varepsilon_0}=P_{\varepsilon_0-\varepsilon}*\nu$ if and only if no zero has $\beta>\frac12+\varepsilon$. The proof runs a Poisson integral of $\nu$ down from $\varepsilon_0$ and uses harmonic continuation: a positive $\nu$ would make $\rho_\varepsilon(\theta)$ harmonic in $(\theta,\varepsilon)$ near a zero with $\beta-\frac12>\varepsilon$, where $\xi'/\xi$ has a pole.

Asking whether the deconvolution \emph{can be done} has two meanings. One is constructive: can $P_{\varepsilon_0}^{-1}\rho_{\varepsilon_0}$ be produced, from data we can compute, without assuming where the zeros are? The other is the theorem itself: can it be shown to be a positive measure? We answer the first completely (Section~\ref{dc:sec:weil}) and reduce the second to a two parameter family of finite problems (Sections~\ref{dc:sec:damped} and~\ref{dc:sec:ladder}), whose behaviour we partly prove and compute.
